
\documentclass[11pt,a4paper]{article}

\usepackage[a4paper,margin=2.2cm]{geometry}
\usepackage{microtype}
\usepackage[T1]{fontenc}
\usepackage{lmodern}
\usepackage{hyperref}
\usepackage{titlesec}
\usepackage{fancyhdr}
\usepackage{enumitem}
\usepackage{longtable}
\usepackage{array}
\usepackage{booktabs}
\usepackage{tabularx}
\usepackage{xcolor}
\usepackage{listings}
\usepackage{framed}
\usepackage{graphicx}
\usepackage{float}
\usepackage{parskip}

\hypersetup{
  colorlinks=true,
  linkcolor=blue!50!black,
  urlcolor=blue!50!black,
  pdftitle={PDF to Markdown Pipeline - Full Project Description},
  pdfauthor={OpenAI}
}

\definecolor{Accent}{HTML}{1F4E79}
\definecolor{Soft}{HTML}{F4F7FA}
\definecolor{CodeBg}{HTML}{F7F7F7}
\definecolor{CodeFrame}{HTML}{CCCCCC}

\titleformat{\section}{\Large\bfseries\color{Accent}}{\thesection.}{0.6em}{}
\titleformat{\subsection}{\large\bfseries\color{Accent}}{\thesubsection}{0.6em}{}
\titleformat{\subsubsection}{\normalsize\bfseries}{\thesubsubsection}{0.6em}{}

\pagestyle{fancy}
\fancyhf{}
\lhead{\textit{PDF to Markdown Pipeline}}
\rhead{\textit{Architecture Specification}}
\cfoot{\thepage}

\lstdefinestyle{custom}{
  basicstyle=\ttfamily\small,
  backgroundcolor=\color{CodeBg},
  frame=single,
  rulecolor=\color{CodeFrame},
  breaklines=true,
  columns=fullflexible,
  keepspaces=true,
  showstringspaces=false
}
\lstset{style=custom}

\newcommand{\docnote}[1]{
\begin{framed}
\noindent\textbf{Note.} #1
\end{framed}
}

\begin{document}

\begin{titlepage}
\centering
\vspace*{3cm}
{\Huge\bfseries PDF to Markdown Pipeline\par}
\vspace{0.5cm}
{\Large Full Project Description, Architecture, Schemes, Tools, and Build Plan\par}
\vspace{1.2cm}
{\large Local-first, Docker-based, artifact-driven document processing platform\par}
\vspace{2cm}

\begin{tabular}{>{\bfseries}p{4.2cm}p{8.6cm}}
Document type: & Technical architecture specification \\
Primary goal: & Build a complete product skeleton that can grow from placeholders into a production system \\
Main input: & PDF document \\
Main outputs: & Markdown, structured JSON, and full intermediate artifacts \\
Runtime model: & Multiple Docker containers, one service per responsibility \\
Communication: & FastAPI over internal Docker network for phase 1 \\
Version target: & Phase-1 / Phase-2 project planning document \\
\end{tabular}

\vfill
{\large \today\par}
\end{titlepage}

\tableofcontents
\newpage

\section{Executive Summary}

This project is a modular platform that receives a PDF document and produces a Markdown representation together with a structured JSON representation and a full set of intermediate artifacts. The system is intentionally designed as a \textbf{pipeline of Docker containers}, not as one large monolithic application.

The key architectural idea is simple:

\begin{itemize}[leftmargin=1.2cm]
\item one service receives and coordinates jobs;
\item one service prepares pages from the PDF;
\item one service detects layout blocks;
\item one service routes blocks to specialized workers;
\item different workers process text, formulas, tables, and images;
\item one service determines reading order;
\item one service assembles final Markdown and JSON;
\item one service validates the final result.
\end{itemize}

This makes the product easier to build, easier to test, easier to debug, and much easier to extend when different OCR or vision models are introduced later.

\section{What the Product Must Do}

The product must eventually support the following end-to-end behavior:

\begin{enumerate}[leftmargin=1.2cm]
\item Receive a PDF from a client.
\item Create a job id and a job folder.
\item Store the original document and job metadata.
\item Convert the PDF into page images.
\item Detect layout blocks on each page.
\item Route blocks to the correct worker.
\item Extract content from each block:
  \begin{itemize}
  \item text blocks;
  \item formulas;
  \item tables;
  \item figures and image-like regions.
  \end{itemize}
\item Reconstruct logical reading order.
\item Assemble a final Markdown file and structured JSON file.
\item Validate results and mark the job as completed, degraded, or failed.
\item Keep all intermediate artifacts so the whole process can be inspected later.
\end{enumerate}

\section{What the First Version Must Actually Do}

The first version does \textbf{not} need perfect OCR or advanced models. The first version needs a \textbf{correct skeleton}.

That means phase 1 should provide:

\begin{itemize}[leftmargin=1.2cm]
\item a clean monorepo;
\item one Dockerfile per service;
\item a working docker-compose setup;
\item a hello-world orchestrator;
\item placeholder services for all major pipeline stages;
\item stable request/response schemas;
\item a job folder layout;
\item a clear Git and model-storage strategy.
\end{itemize}

The point is to prove the \textbf{architecture}, not model accuracy.

\section{Core Architectural Principles}

\subsection{One Service = One Responsibility}

Each service should do one thing well.

Examples:

\begin{itemize}[leftmargin=1.2cm]
\item the layout service finds blocks;
\item the text OCR service reads text regions;
\item the formula OCR service reads formulas;
\item the table OCR service reads tables;
\item the image classifier labels figure-like regions;
\item the orchestrator coordinates the pipeline.
\end{itemize}

This prevents the project from collapsing into one huge service with mixed responsibilities.

\subsection{One Service = One Docker Container}

Each service runs inside its own container and therefore has its own:

\begin{itemize}[leftmargin=1.2cm]
\item Python version;
\item packages;
\item system libraries;
\item CUDA-related dependencies if needed;
\item file-level runtime environment.
\end{itemize}

This is crucial because different workers may use different frameworks such as Paddle, PyTorch, custom OCR stacks, or different local LLMs.

\subsection{Stable Contracts Between Services}

Services must communicate using stable schemas. A worker can change its internal model later, but its API contract should stay compatible.

This allows:

\begin{itemize}[leftmargin=1.2cm]
\item easy model replacement;
\item cleaner testing;
\item easier integration of alternatives;
\item less coupling between services.
\end{itemize}

\subsection{Artifact-Driven Design}

Every pipeline stage should write artifacts to disk. This means the system never hides all intermediate state inside memory only.

This gives:

\begin{itemize}[leftmargin=1.2cm]
\item reproducibility;
\item step-by-step debugging;
\item partial reruns;
\item easier quality inspection;
\item future benchmark datasets;
\item operational transparency.
\end{itemize}

\section{High-Level System Scheme}

\begin{lstlisting}
Client
  |
  v
Orchestrator
  |
  v
Page Preparation
  |
  v
Layout Master
  |
  v
Block Router
  |------> Text OCR
  |------> Formula OCR
  |------> Table OCR
  |------> Image Classifier
                  |
                  v
             Image OCR / Figure Worker
  |
  v
Reading Order
  |
  v
Assembly
  |
  v
Validation
  |
  v
document.md + document.json + debug artifacts + logs
\end{lstlisting}

\section{Detailed Description of Each Service}

\subsection{1. Orchestrator}

\textbf{Role:} central control plane for a job.

\textbf{Responsibilities:}
\begin{itemize}[leftmargin=1.2cm]
\item expose public API for upload and status;
\item create job ids;
\item create job directories;
\item call all services in the right order;
\item write job metadata and trace files;
\item collect service outputs;
\item update stage statuses;
\item decide final job state.
\end{itemize}

\textbf{Must not do:}
\begin{itemize}[leftmargin=1.2cm]
\item real OCR;
\item model inference;
\item table extraction logic;
\item formula recognition;
\item figure classification.
\end{itemize}

\textbf{Reasoning:} the orchestrator should stay focused on pipeline control, not on document intelligence.

\subsection{2. Page Preparation Service}

\textbf{Role:} normalize PDF input into page-level images.

\textbf{Responsibilities:}
\begin{itemize}[leftmargin=1.2cm]
\item store original file;
\item render PDF pages to PNG or another page-image format;
\item create page numbering;
\item optionally perform preprocessing:
  \begin{itemize}
  \item deskew;
  \item denoise;
  \item contrast normalization;
  \item resizing.
  \end{itemize}
\item write per-page metadata.
\end{itemize}

\textbf{Output examples:}
\begin{itemize}[leftmargin=1.2cm]
\item \texttt{assets/pages/page\_0001.png}
\item \texttt{assets/pages/page\_0002.png}
\item \texttt{meta.json}
\end{itemize}

\subsection{3. Layout Master}

\textbf{Role:} detect logical layout blocks on pages.

\textbf{Typical block types:}
\begin{itemize}[leftmargin=1.2cm]
\item title;
\item text;
\item list;
\item caption;
\item table;
\item formula;
\item figure / image;
\item header;
\item footer;
\item footnote;
\item reference.
\end{itemize}

\textbf{Responsibilities:}
\begin{itemize}[leftmargin=1.2cm]
\item run layout detection model;
\item output block coordinates and types;
\item save raw model output;
\item save normalized layout output;
\item save debug overlays with boxes and labels.
\end{itemize}

\textbf{Important rule:} layout service says \emph{what exists and where}. It does not read text and does not assemble Markdown.

\subsection{4. Block Router}

\textbf{Role:} crop detected blocks and dispatch them to the correct worker.

\textbf{Responsibilities:}
\begin{itemize}[leftmargin=1.2cm]
\item read normalized layout results;
\item crop each block from its page image;
\item decide routing target;
\item call specialized workers;
\item collect worker outputs and save them.
\end{itemize}

\textbf{Routing examples:}
\begin{itemize}[leftmargin=1.2cm]
\item text-like block $\rightarrow$ text OCR;
\item formula block $\rightarrow$ formula OCR;
\item table block $\rightarrow$ table OCR;
\item image-like block $\rightarrow$ image classifier.
\end{itemize}

\subsection{5. Text OCR}

\textbf{Role:} process text-like blocks.

\textbf{Typical inputs:}
\begin{itemize}[leftmargin=1.2cm]
\item paragraphs;
\item titles;
\item captions;
\item list items;
\item references;
\item headers and footers.
\end{itemize}

\textbf{Typical output fields:}
\begin{itemize}[leftmargin=1.2cm]
\item block id;
\item recognized text;
\item confidence;
\item optional formatting hints;
\item raw model output.
\end{itemize}

\subsection{6. Formula OCR}

\textbf{Role:} process formula blocks and convert them into a useful format such as LaTeX.

\textbf{Typical output fields:}
\begin{itemize}[leftmargin=1.2cm]
\item block id;
\item formula LaTeX;
\item confidence;
\item raw model response.
\end{itemize}

\textbf{Why separate:} formulas need different recognition logic than plain OCR text.

\subsection{7. Table OCR}

\textbf{Role:} turn table crops into structured table data.

\textbf{Typical output fields:}
\begin{itemize}[leftmargin=1.2cm]
\item block id;
\item rows and columns;
\item cell coordinates if supported;
\item table title if detected;
\item markdown-ready representation;
\item JSON-ready representation.
\end{itemize}

\textbf{Why separate:} tables require structure understanding, not only text recognition.

\subsection{8. Image Classifier}

\textbf{Role:} classify image-like blocks before deeper handling.

\textbf{Useful labels:}
\begin{itemize}[leftmargin=1.2cm]
\item chart;
\item diagram;
\item photo;
\item illustration;
\item scanned-text-image;
\item other.
\end{itemize}

\textbf{Why this service matters:} different image categories need different downstream treatment.

\subsection{9. Image OCR / Figure Worker}

\textbf{Role:} handle image-like blocks after classification.

\textbf{Possible tasks:}
\begin{itemize}[leftmargin=1.2cm]
\item OCR inside figures;
\item placeholder figure description;
\item scientific figure understanding later;
\item chart or diagram forwarding later.
\end{itemize}

\subsection{10. Reading Order Service}

\textbf{Role:} reconstruct natural human reading order.

\textbf{Why needed:}
\begin{itemize}[leftmargin=1.2cm]
\item multi-column layouts;
\item captions under images;
\item sidebars;
\item header/footer suppression;
\item complex block sequences.
\end{itemize}

\subsection{11. Assembly Service}

\textbf{Role:} build final document representations.

\textbf{Responsibilities:}
\begin{itemize}[leftmargin=1.2cm]
\item merge ordered block outputs;
\item generate \texttt{document.md};
\item generate \texttt{document.json};
\item preserve hierarchy where possible;
\item serialize formulas and tables correctly.
\end{itemize}

\subsection{12. Validation Service}

\textbf{Role:} verify the final output is complete and coherent.

\textbf{Responsibilities:}
\begin{itemize}[leftmargin=1.2cm]
\item verify all expected artifacts exist;
\item validate schema structure;
\item detect suspicious empty blocks;
\item mark job status as completed, degraded, or failed.
\end{itemize}

\section{Communication Between Containers}

For phase 1, services should communicate through \textbf{HTTP on the Docker network}. FastAPI is a very good fit because it gives:

\begin{itemize}[leftmargin=1.2cm]
\item fast development;
\item clean REST endpoints;
\item automatic request/response validation;
\item interactive docs for testing;
\item simple JSON contracts.
\end{itemize}

\subsection{Important Concept}

Containers do not communicate because of FastAPI itself. They communicate because they share a Docker network and use HTTP. FastAPI is simply the web framework that exposes the HTTP endpoints.

\subsection{Internal Addressing}

Inside Docker Compose, services can call each other by service name:

\begin{lstlisting}
http://orchestrator:8000
http://layout_master:8000
http://text_ocr:8000
http://formula_ocr:8000
http://table_ocr:8000
http://image_classifier:8000
\end{lstlisting}

\subsection{Recommended Minimal Endpoints Per Worker}

\begin{lstlisting}
GET  /health
GET  /capabilities
POST /process
\end{lstlisting}

This keeps every service predictable.

\section{Project Repository Structure}

A recommended monorepo layout is:

\begin{lstlisting}
pdf-md-pipeline/
├── .gitignore
├── .dockerignore
├── .env.example
├── README.md
├── LICENSE
├── docker-compose.yml
├── docs/
│   ├── ARCHITECTURE.md
│   ├── ARTIFACT_CONTRACT.md
│   ├── API_CONTRACTS.md
│   ├── DEVELOPMENT.md
│   └── DEPLOYMENT.md
├── configs/
│   ├── orchestrator/
│   ├── page_preparation/
│   ├── layout_master/
│   ├── block_router/
│   ├── text_ocr/
│   ├── formula_ocr/
│   ├── table_ocr/
│   ├── image_classifier/
│   ├── image_ocr/
│   ├── reading_order/
│   ├── assembly/
│   └── validation/
├── shared/
│   ├── schemas/
│   ├── clients/
│   └── utils/
├── services/
│   ├── orchestrator/
│   ├── page_preparation/
│   ├── layout_master/
│   ├── block_router/
│   ├── text_ocr/
│   ├── formula_ocr/
│   ├── table_ocr/
│   ├── image_classifier/
│   ├── image_ocr/
│   ├── reading_order/
│   ├── assembly/
│   └── validation/
├── scripts/
├── data/
│   ├── jobs/
│   ├── models/
│   └── cache/
├── examples/
└── tests/
\end{lstlisting}

\section{Service Internal Structure}

Each service should have a similar internal skeleton to reduce friction:

\begin{lstlisting}
services/text_ocr/
├── app/
│   ├── main.py
│   ├── api/
│   ├── core/
│   ├── models/
│   └── config.py
├── tests/
├── requirements.txt
└── Dockerfile
\end{lstlisting}

This consistency matters. It makes the whole monorepo easier to navigate.

\section{Artifact Contract}

Every job should have a predictable folder structure. Example:

\begin{lstlisting}
data/jobs/job_000001/
├── original/
│   └── input.pdf
├── preprocessed/
├── assets/
│   ├── pages/
│   │   ├── page_0001.png
│   │   └── page_0002.png
│   ├── layout/
│   │   ├── page_0001.with_boxes.png
│   │   └── page_0001.with_boxes_labeled.png
│   ├── text_blocks/
│   ├── tables/
│   ├── formulas/
│   └── images/
├── debug/
│   ├── layout_raw.json
│   ├── layout_normalized.json
│   └── service_responses/
├── meta.json
├── pipeline_trace.json
├── logs.jsonl
├── document.json
└── document.md
\end{lstlisting}

\subsection{Why This Contract Is Important}

The artifact contract creates a stable operational truth. Anyone can inspect a job and understand:

\begin{itemize}[leftmargin=1.2cm]
\item what the input was;
\item what each stage produced;
\item where a stage failed;
\item whether the final result is trustworthy.
\end{itemize}

\section{Job State Model}

Every job should keep an explicit state. Suggested stage states:

\begin{itemize}[leftmargin=1.2cm]
\item pending;
\item running;
\item completed;
\item failed;
\item degraded;
\item skipped.
\end{itemize}

Suggested high-level job lifecycle:

\begin{lstlisting}
created
uploaded
preparing_pages
layout_running
layout_done
routing_blocks
text_ocr_running
formula_ocr_running
table_ocr_running
image_classification_running
reading_order_done
assembly_done
validation_done
completed
\end{lstlisting}

Possible end states:
\begin{itemize}[leftmargin=1.2cm]
\item completed;
\item degraded;
\item failed.
\end{itemize}

\section{Docker and Configuration Model}

\subsection{Why We Still Need Environment Variables}

Docker isolates runtime, but it does not remove configuration. Services still need to know:

\begin{itemize}[leftmargin=1.2cm]
\item where other services live;
\item what port to listen on;
\item where models are mounted;
\item where job artifacts are written;
\item CPU or GPU mode;
\item timeout and batch settings;
\item log level.
\end{itemize}

\subsection{Recommended Files}

\begin{itemize}[leftmargin=1.2cm]
\item \texttt{.env.example} goes into Git;
\item real \texttt{.env} stays local and is not committed.
\end{itemize}

\subsection{Example .env Values}

\begin{lstlisting}
ORCHESTRATOR_PORT=8000
ARTIFACT_ROOT=/data/jobs
MODEL_ROOT=/models
LAYOUT_MASTER_URL=http://layout_master:8000
TEXT_OCR_URL=http://text_ocr:8000
FORMULA_OCR_URL=http://formula_ocr:8000
TABLE_OCR_URL=http://table_ocr:8000
IMAGE_CLASSIFIER_URL=http://image_classifier:8000
IMAGE_OCR_URL=http://image_ocr:8000
LOG_LEVEL=INFO
\end{lstlisting}

\section{Docker Compose Strategy}

For the first version, only the orchestrator should usually expose a host port. Worker services can stay internal.

\begin{lstlisting}
Client browser --> localhost:8000 --> orchestrator
orchestrator --> internal docker network --> workers
\end{lstlisting}

That keeps the local deployment cleaner and avoids unnecessary host-level port clutter.

\docnote{As the system grows, different services may need different base images, Python versions, GPU access settings, and volume mounts. This is exactly why one container per service is the right foundation.}

\section{Python, Virtual Environments, and Containers}

A very important conceptual point from the project discussion is this:

\begin{itemize}[leftmargin=1.2cm]
\item a local \texttt{.venv} is optional and mainly useful for development convenience;
\item the real runtime environment is each Docker container;
\item each container already has its own Python interpreter and its own installed packages.
\end{itemize}

So the correct mental model is not ``one project, one venv''.

The correct model is:

\begin{lstlisting}
orchestrator container -> its own Python + deps
layout container       -> its own Python + deps
text OCR container     -> its own Python + deps
formula container      -> its own Python + deps
\end{lstlisting}

This is what allows different stacks to coexist.

\section{Git and GitHub Strategy}

\subsection{What Must Be Versioned}

GitHub should store the durable project definition:

\begin{itemize}[leftmargin=1.2cm]
\item source code;
\item Dockerfiles;
\item docker-compose;
\item configs;
\item docs;
\item schemas;
\item scripts;
\item tiny sample files;
\item CI files later;
\item \texttt{.env.example}.
\end{itemize}

\subsection{What Must Not Be Versioned}

GitHub should \textbf{not} store:

\begin{itemize}[leftmargin=1.2cm]
\item huge model weights;
\item generated artifacts;
\item runtime logs;
\item local caches;
\item local \texttt{.env};
\item \texttt{.venv};
\item page images produced during jobs;
\item temporary files.
\end{itemize}

\subsection{Recommended .gitignore Sections}

\begin{lstlisting}
# Python
__pycache__/
*.pyc
.venv/
.pytest_cache/
.mypy_cache/
.ruff_cache/

# Runtime data
data/jobs/
data/cache/
logs/
tmp/
output/

# Models
models/
data/models/
*.safetensors
*.bin
*.pt
*.pth
*.onnx
*.pdmodel
*.pdiparams
*.ckpt

# IDE / OS
.vscode/
.idea/
.DS_Store
Thumbs.db

# Secrets
.env
\end{lstlisting}

\section{Large Models and Multi-Machine Work}

A critical operational question in this project is how to work across a laptop, a home PC, or other machines when the system may use huge local LLMs or OCR models.

\subsection{What GitHub Is For}

GitHub is for:
\begin{itemize}[leftmargin=1.2cm]
\item code;
\item project structure;
\item Docker definitions;
\item scripts;
\item configs;
\item documentation.
\end{itemize}

\subsection{What GitHub Is Not For}

GitHub is not the right storage layer for hundreds of gigabytes of local model weights. Large local models should live outside the repository.

\subsection{Correct Strategy}

Use this split:

\begin{itemize}[leftmargin=1.2cm]
\item \textbf{GitHub:} code and architecture;
\item \textbf{local disks / NAS / model storage:} weights and caches;
\item \textbf{Docker mounts:} runtime access to those models.
\end{itemize}

\subsection{Per-Machine Example}

Laptop:
\begin{lstlisting}
D:\models-lite\
\end{lstlisting}

Home PC:
\begin{lstlisting}
C:\models\
\end{lstlisting}

Linux workstation:
\begin{lstlisting}
/srv/models/
\end{lstlisting}

Each machine clones the same repository, but the local \texttt{.env} points containers to that machine's model directory.

\subsection{Recommended Operational Pattern}

\begin{enumerate}[leftmargin=1.2cm]
\item clone the same GitHub repo on each machine;
\item create local \texttt{.env};
\item define machine-specific model paths;
\item mount model directories into containers;
\item run docker-compose locally.
\end{enumerate}

\subsection{Model Setup Scripts}

The repository should include scripts such as:

\begin{itemize}[leftmargin=1.2cm]
\item \texttt{scripts/fetch-models.sh}
\item \texttt{scripts/fetch-models.ps1}
\end{itemize}

These scripts can:

\begin{itemize}[leftmargin=1.2cm]
\item download approved models from their official source;
\item copy models from a NAS or shared location;
\item verify checksums;
\item place files into expected folders.
\end{itemize}

That is far better than trying to store giant binary weights in Git.

\section{Enumeration of Main Instruments and Technologies}

This section lists the main instruments the project will use or is designed to support.

\subsection{Infrastructure and Packaging}

\begin{itemize}[leftmargin=1.2cm]
\item \textbf{Git} -- local version control.
\item \textbf{GitHub} -- remote code repository and collaboration hub.
\item \textbf{Docker} -- runtime isolation per service.
\item \textbf{Docker Compose} -- local multi-service orchestration.
\end{itemize}

\subsection{Programming and API Layer}

\begin{itemize}[leftmargin=1.2cm]
\item \textbf{Python} -- primary implementation language.
\item \textbf{FastAPI} -- service HTTP APIs.
\item \textbf{Pydantic} -- schema validation.
\item \textbf{Uvicorn} -- ASGI server for services.
\end{itemize}

\subsection{Shared Engineering Tools}

\begin{itemize}[leftmargin=1.2cm]
\item \textbf{pytest} -- tests.
\item \textbf{ruff / mypy} -- code quality and typing later.
\item \textbf{logging / json logs} -- observability.
\end{itemize}

\subsection{Possible Document / Vision Engines}

This architecture is intentionally engine-agnostic. Different workers can later use different stacks, for example:

\begin{itemize}[leftmargin=1.2cm]
\item layout model service;
\item OCR model service;
\item formula-recognition model service;
\item table extraction service;
\item local vision-language model for image classification.
\end{itemize}

The architecture cares more about \textbf{contracts} than about locking the project to one model.

\section{Suggested API Contracts}

\subsection{Health}

\begin{lstlisting}
GET /health

Response:
{
  "status": "ok"
}
\end{lstlisting}

\subsection{Capabilities}

\begin{lstlisting}
GET /capabilities

Response:
{
  "service": "text_ocr",
  "version": "0.1.0",
  "mode": "placeholder",
  "supported_inputs": ["image_crop"]
}
\end{lstlisting}

\subsection{Process}

\begin{lstlisting}
POST /process

Request:
{
  "job_id": "job_000001",
  "block_id": "p1_b7",
  "input_path": "/data/jobs/job_000001/assets/text_blocks/p1_b7.png"
}
\end{lstlisting}

\begin{lstlisting}
Response:
{
  "job_id": "job_000001",
  "block_id": "p1_b7",
  "status": "completed",
  "output_path": "/data/jobs/job_000001/assets/text_blocks/p1_b7.result.json"
}
\end{lstlisting}

\section{Suggested Shared Schemas}

The \texttt{shared/schemas} package should define canonical structures such as:

\begin{itemize}[leftmargin=1.2cm]
\item \texttt{Job}
\item \texttt{StageStatus}
\item \texttt{LayoutBlock}
\item \texttt{OCRResult}
\item \texttt{FormulaResult}
\item \texttt{TableResult}
\item \texttt{ImageClassificationResult}
\item \texttt{ArtifactRecord}
\end{itemize}

This keeps all services speaking the same language.

\section{Ready Build Plan}

This section is the practical plan you can follow almost directly while building the project.

\subsection{Stage 0: Create the Repository}

\begin{enumerate}[leftmargin=1.2cm]
\item Create a new Git repository.
\item Create the monorepo folder tree.
\item Add \texttt{README.md}, \texttt{LICENSE}, \texttt{.gitignore}, \texttt{.dockerignore}, and \texttt{.env.example}.
\item Commit the empty structure first.
\end{enumerate}

\subsection{Stage 1: Create Placeholder Services}

For each planned service:

\begin{enumerate}[leftmargin=1.2cm]
\item create \texttt{app/main.py};
\item add \texttt{requirements.txt};
\item add a service-specific \texttt{Dockerfile};
\item implement:
  \begin{itemize}
  \item \texttt{/health}
  \item \texttt{/capabilities}
  \item \texttt{/process}
  \end{itemize}
\item return placeholder JSON from \texttt{/process}.
\end{enumerate}

\subsection{Stage 2: Create the Orchestrator}

\begin{enumerate}[leftmargin=1.2cm]
\item create upload endpoint;
\item generate \texttt{job\_id};
\item create the artifact folder tree;
\item call page preparation service;
\item call layout service;
\item for now call placeholder workers;
\item write a placeholder \texttt{document.md};
\item expose job status endpoint.
\end{enumerate}

\subsection{Stage 3: Add docker-compose}

\begin{enumerate}[leftmargin=1.2cm]
\item define one service per container;
\item expose only orchestrator to host;
\item mount:
  \begin{itemize}
  \item artifact root;
  \item model directories;
  \item optional config volumes.
  \end{itemize}
\item verify orchestrator can reach workers by service name.
\end{enumerate}

\subsection{Stage 4: Add Shared Schemas}

\begin{enumerate}[leftmargin=1.2cm]
\item define Pydantic models for job, stage, block, and result types;
\item use these schemas in request and response bodies;
\item centralize status enums and artifact metadata.
\end{enumerate}

\subsection{Stage 5: Add Basic Integration Tests}

\begin{enumerate}[leftmargin=1.2cm]
\item health checks for all services;
\item orchestrator pipeline smoke test;
\item verify placeholder output files exist;
\item verify \texttt{document.md} and \texttt{document.json} are generated.
\end{enumerate}

\subsection{Stage 6: Replace Placeholders One by One}

\begin{enumerate}[leftmargin=1.2cm]
\item first replace page preparation with real PDF rendering;
\item then replace layout;
\item then text OCR;
\item then formula OCR;
\item then table OCR;
\item then image classification;
\item then reading order and validation improvements.
\end{enumerate}

\subsection{Stage 7: Production Hardening Later}

\begin{enumerate}[leftmargin=1.2cm]
\item retries;
\item degraded-mode logic;
\item queue-based execution if needed;
\item better metrics;
\item distributed workers;
\item web UI or admin panel later.
\end{enumerate}

\section{Recommended Reading Order for Building the Project}

If you want to study the project and build it only from this plan, use this order:

\begin{enumerate}[leftmargin=1.2cm]
\item read the high-level architecture section;
\item read the service descriptions;
\item read the repository structure;
\item read the artifact contract;
\item read the Docker and configuration model;
\item read the Git and model-storage strategy;
\item follow the ready build plan stage by stage;
\item only then integrate real models.
\end{enumerate}

That order will keep the project understandable.

\section{Final Architecture Summary}

The project should be understood as this:

\begin{itemize}[leftmargin=1.2cm]
\item \textbf{GitHub repository} = code, docs, configs, Dockerfiles, schemas.
\item \textbf{Docker Compose} = starts the whole platform locally.
\item \textbf{Each service} = one container, one responsibility, one runtime.
\item \textbf{FastAPI} = simple HTTP communication layer.
\item \textbf{Artifact folders} = persistent operational truth of each job.
\item \textbf{Local model folders} = heavy weights outside Git.
\end{itemize}

This design is the right base for a system that starts as a clean placeholder scaffold and later grows into a serious offline PDF-to-Markdown platform.

\section{Concise Master Scheme}

\begin{lstlisting}
               +----------------------+
               |       Client         |
               +----------+-----------+
                          |
                          v
               +----------------------+
               |    Orchestrator      |
               | job control + API    |
               +----------+-----------+
                          |
                          v
               +----------------------+
               |  Page Preparation    |
               +----------+-----------+
                          |
                          v
               +----------------------+
               |    Layout Master     |
               +----------+-----------+
                          |
                          v
               +----------------------+
               |     Block Router     |
               +---+----+----+----+---+
                   |    |    |    |
                   v    v    v    v
               +------+ +------+ +------+
               |Text  | |Form. | |Table |
               | OCR  | | OCR  | | OCR  |
               +------+ +------+ +------+
                          |
                          v
                    +-----------+
                    |  Image    |
                    |Classifier |
                    +-----+-----+
                          |
                          v
                    +-----------+
                    | Image OCR |
                    +-----+-----+
                          |
                          v
                    +-----------+
                    | Reading   |
                    |  Order    |
                    +-----+-----+
                          |
                          v
                    +-----------+
                    | Assembly  |
                    +-----+-----+
                          |
                          v
                    +-----------+
                    |Validation |
                    +-----+-----+
                          |
                          v
               +----------------------+
               | Markdown + JSON +    |
               | artifacts + logs     |
               +----------------------+
\end{lstlisting}

\end{document}
